%%%%%%%%%%%%%%%%%%%%%%%%%%%%%%%%%%%%%%%%%
% TEMPLATE: Academic CV (cv.cls)
% FIXED sections contain actual content — copy verbatim during generation
% JD-DEPENDENT sections have [GENERATE: ...] placeholders — replace per JD
%
% SETUP INSTRUCTIONS:
% 1. Fill [CONFIG: ...] with values from your config.md
% 2. Fill ALL FIXED sections with your actual content
% 3. Leave [GENERATE: ...] markers -- Claude fills these per JD
% 4. After setup, compile and calibrate the line budget below
%
% PAGE 1 RULE: First bullet of first experience MUST be 2L (not 3L).
% A 3L first bullet overflows page 1. If top achievement needs 3L,
% make it bullet #2 and lead with a strong 2L bullet.
%%%%%%%%%%%%%%%%%%%%%%%%%%%%%%%%%%%%%%%%%

\documentclass[letterpaper]{cv} % Notion-style; drop [letterpaper] for A4
\usepackage{fontawesome}
\usepackage{graphicx}
\hypersetup{
    pdftitle = {AZ S. Zhang - Academic CV},
    pdfauthor = {AZ S. Zhang}
}
% Styling (Inter font, Notion colors, geometry, links, footer, mhchem) is set in cv.cls.
% Do NOT load geometry, fancyhdr, or hyperref here -- it will clash with the class.

%----------------------------------------------------------------------------------------
%	HEADER — FIXED
%----------------------------------------------------------------------------------------
\name{AZ S. Zhang}
\address{\faEnvelope~\href{mailto:azz@stanford.edu}{azz@stanford.edu} \\ \faGlobe~\href{https://azzhang.dev/}{azzhang.dev} \\ \faGithub~\href{https://github.com/akmazian}{GitHub}}
\address{\faPhone~+1 (949) 237-3629 \\ \faMapMarker~Stanford, CA \\ \faLinkedin~\href{https://www.linkedin.com/in/az-zhang/}{LinkedIn} \\ \includegraphics[height=1em]{orcid.png}~\href{https://orcid.org/0009-0009-2503-1800}{ORCID: 0009-0009-2503-1800} \\ \includegraphics[height=1em]{GS.png}~\href{https://scholar.google.com/citations?user=dwLg1YkAAAAJ&hl=en}{Google Scholar}}

\begin{document}

%========================================================================================
%  CV LINE BUDGET — calibrate after filling FIXED sections
%========================================================================================
%
% HOW TO CALIBRATE:
% 1. Fill all FIXED sections with your actual content
% 2. Compile with dummy GENERATE content
% 3. Count rendered lines per page
% 4. Update the budget numbers below
%
% TARGET: page count from config.md (Letter, 10pt Inter cv.cls). Do not pad.
% Budget: ~39 one-line / ~23 two-line bullets per page.
% NO TECHNICAL EXPERTISE and NO RESEARCH SUMMARY SECTION (user preference, see config.md).
% 1-2 lines slack at bottom of last page is acceptable.
%
% FIXED = content that never changes across JDs
% JD-DEPENDENT = Claude generates per JD, but line COUNT is fixed
%
% CHARACTER LIMITS (rendered chars at 10pt Inter, ragged-right, cv.cls geometry):
%   Chars per line (CPL): ~92-97
%
%   CV-2L bullet: target 168-182 chars | HARD MAX 190 | orphan >= 65 chars
%   CV-3L bullet: target 250-268 chars | HARD MAX 280 | orphan >= 65 chars
%
%   AIM FOR TARGET MIDDLE (175 for 2L, 260 for 3L) — NOT the hard max.
%   Em-dash (---) counts as 1 char but renders ~2x wide. Budget 2 extra per em-dash.
%   Bold in bullets: ~0.33 penalty/bold char (10 bold chars ≈ lose 3 chars capacity).
%
% EXPERIENCE BULLET BUDGET:
%   After calibrating, determine how many rendered lines are available for
%   experience bullets. Typical: ~45 rendered lines.
%   Mix: 2L (2 rendered lines) + 3L (3 rendered lines)
%   Example: 18x2L + 3x3L = 21 bullets in 45 lines
%========================================================================================

%========================================================================================
%  EDUCATION — FIXED
%========================================================================================

\begin{rSection}{Education}
\textbf{Ph.D. Program in Genetics}\notionsep Sep 2026 -- Present\\
Stanford University, Stanford, CA

\textbf{B.A., Molecular and Cell Biology with Honors}\notionsep Aug 2022 -- Dec 2024\\
University of California, Berkeley, Berkeley, CA\notionsep GPA 3.98/4.00\\
{\color{notiongray}Relevant coursework: General Genetics, Advanced Genetic Analysis, Computational Functional Genomics, Machine Learning, Genetic Design Automation, Biophysical Chemistry, Human Genomics}
\end{rSection}

%========================================================================================
%  RESEARCH EXPERIENCE — MIXED FIXED/JD-DEPENDENT
%  Position headers: FIXED | Sub-theme count: FIXED per position | Content: JD-DEPENDENT
%========================================================================================
% cv.cls rSubsection: {Title}{Dates}{Location}{Mentors/Advisors}
% CRITICAL: cv.cls puts Args 3 and 4 on SEPARATE muted lines (unlike resume.cls!)
% Sub-themes via \subtheme{...}
% NEVER use code folder names as package names.

\begin{rSection}{Research Experience}

% --- Position 1: Weissman Lab. Replace sub-theme names and bullets per JD. Keep sub-theme COUNT fixed.
\begin{rSubsection}{Research Technician}{Jan 2025 -- May 2026}{Weissman Lab, Whitehead Institute for Biomedical Research}{PI: Jonathan S. Weissman}

\subtheme{[GENERATE: Sub-theme 1 name]}
\item [GENERATE: Bullet]
\item [GENERATE: Bullet]
\subtheme{[GENERATE: Sub-theme 2 name]}
\item [GENERATE: Bullet]

\end{rSubsection}

% --- Position 2: Ünal Lab ---
\begin{rSubsection}{Undergraduate Research Assistant}{Spring 2023 -- Fall 2024}{\"{U}nal Lab, Department of Molecular and Cell Biology, UC Berkeley}{PI: El\c{c}in \"{U}nal}

\subtheme{[GENERATE: Sub-theme 1 name]}
\item [GENERATE: Bullet]
\item [GENERATE: Bullet]

\end{rSubsection}

\end{rSection}

%========================================================================================
%  FELLOWSHIPS & HONORS — FIXED
%========================================================================================

\begin{rSection2}{Fellowships \& Honors}
\item \textbf{Stanford Graduate Fellowship}, Stanford University (2026--2029).
\item \textbf{High Distinction in General Scholarship}, UC Berkeley (2024).
\item \textbf{Rose Hills Summer Research Fellowship}, UC Berkeley and the Rose Hills Foundation (2024). Independent experimental design and analysis in the \"{U}nal Lab.
\end{rSection2}

%========================================================================================
%  PUBLICATIONS — FIXED (preprint; see config.md Provenance Flags)
%========================================================================================
% Rule 4 et al. format. NOT peer-reviewed -- never list under "Published".

\begin{rSection}{Publications}
{\color{notiongray}$^\dagger$ equal contribution}

\textbf{Preprints}
\begin{enumerate}[leftmargin=1.5em, labelsep=0.5em, itemsep=0.1em]
\item[1.] Chen-Neumann, Y.~H., Grauvogel, L.$^\dagger$, \textbf{Zhang, A.~S.}$^\dagger$, \textit{et al.} ``The diverse functions of viral microproteins.'' \textit{bioRxiv} (2026). \href{https://doi.org/10.64898/2026.08.23.746553}{doi:10.64898/2026.08.23.746553}
\end{enumerate}
\end{rSection}

%========================================================================================
%  POSTERS & PRESENTATIONS — FIXED
%========================================================================================

\begin{rSection2}{Posters \& Presentations}
\item ``Functional annotation of the viral dark proteome.'' \textit{MIT Biology IAP Seminar Poster Session}, Cambridge, MA (Jan 2026). Poster.
\item ``High-throughput functional genomics of pan-viral microproteomes reveals novel viral oncogenes.'' \textit{Whitehead Institute Annual Scientific Retreat} (Sep 2025). Poster.
\item ``Optimizing reverse chromatin immunoprecipitation using dCas9 for locus-specific proteomics.'' \textit{UC Berkeley MCB Honors Thesis Poster Session}, Berkeley, CA (Dec 2024). Poster.
\end{rSection2}

%========================================================================================
%  TEACHING — FIXED
%========================================================================================

\begin{rSection2}{Teaching}
\item \textbf{Undergraduate Teaching Assistant}, Computational Functional Genomics, UC Berkeley (Fall 2024).
\item \textbf{Undergraduate Teaching Assistant}, Data Science for Biology, UC Berkeley (Spring 2024).
\end{rSection2}

\end{document}
